\section{微分の公式}\label{節:微分法則}
基本的な関数を一つずつ定義から微分しても,複雑な式を毎回展開するのは大変である.そこで,式の作られ方と微分の関係を先に調べる.

\begin{theorem}{和・積・商の微分}
$f,g$ が微分可能で,\,$c$ が定数なら,
\begin{align*}
 (f+g)'&=f'+g',& (cf)'&=cf',\\
 (fg)'&=f'g+fg',&
 \left(\frac fg\right)'&=\frac{f'g-fg'}{g^2}\quad(g\ne0)
\end{align*}
が成り立つ.
\end{theorem}

和と定数倍は差商を分ければ分かる.積の公式では,二つの因子が同時に変わることに注意する.途中の項を一つ足して引けば,
\begin{align*}
 &f(x+h)g(x+h)-f(x)g(x)\\
 &\quad=\{f(x+h)-f(x)\}g(x+h)
       +f(x)\{g(x+h)-g(x)\}.
\end{align*}
これを $h$ で割って極限をとる.微分可能な $g$ は連続だから $g(x+h)\to g(x)$ となり,積の公式が得られる.積を微分して単に $f'g'$ としてはいけない.

逆数については,
\[
 \cfrac{\cfrac1{g(x+h)}-\cfrac1{g(x)}}h
 =-\frac{g(x+h)-g(x)}{h\,g(x+h)g(x)}
 \longrightarrow-\frac{g'(x)}{g(x)^2}.
\]
$g(x)\ne0$ と連続性により,十分近くでは $g(x+h)\ne0$ となる.この逆数の公式と積の公式から,商の公式も導ける.

例えば,$x^2(x+1)$では外側が積なので,二つの因子を順に微分して
\[
 \{x^2(x+1)\}'=2x(x+1)+x^2
\]
となる.一方,$\frac{x}{x+1}$には商の公式を使い,
\[
 \left(\frac{x}{x+1}\right)'
 =\frac{1\cdot(x+1)-x\cdot1}{(x+1)^2}
 =\frac1{(x+1)^2}\quad(x\ne-1)
\]
と求められる.

\begin{theorem}{合成関数の微分}\label{定理:合成微分}
$u=g(x)$ が $x$ で微分可能で,\,$f$ が $u=g(x)$ で微分可能なら,
\[
 \frac{\dd}{\dd x}f(g(x))=f'(g(x))g'(x)
\]
である.
\end{theorem}

外側の関数の変化率に,内側の変数の変化率を掛ける公式である.証明のため,\,$\Delta u=g(x+h)-g(x)$ とおく.\,$g$ の連続性から $\Delta u\to0$ であり,\,$f$ の微分可能性から
\[
 f(u+\Delta u)-f(u)=f'(u)\Delta u+\Delta u\,\eta(\Delta u),
 \qquad \eta(\Delta u)\to0
\]
と書ける.\,$\eta(0)=0$ と定めれば,\,$\Delta u=0$ の場合もこの等式は使える.\,$h$ で割ると
\[
 \frac{f(g(x+h))-f(g(x))}h
 =\{f'(g(x))+\eta(\Delta u)\}
   \frac{g(x+h)-g(x)}h.
\]
右辺の極限は $f'(g(x))g'(x)$ になる.途中で $\Delta u$ を分母に置いていないので,内側の関数が同じ値を取る場合も含めて証明できた.

例えば$(x^2+1)^2$は,内側を$u=x^2+1$,外側を$u^2$と見る.外側を微分した$2u$に,内側を微分した$2x$を掛けて
\[
 \{(x^2+1)^2\}'=2(x^2+1)\cdot2x=4x(x^2+1)
\]
となる.内側の$x^2+1$は残したまま,その微分を掛けている.
